\part{Beispiel}
\chapterimage{Bilder_und_Co/analyse.jpg} % Chapter heading image
%----------------------------------------------------------------------------------------
% --- KAPITEL: ANFORDERUNGEN UND SZENARIEN ---
\chapter{Radarauswahl \& Wetterklassifikation}\label{chap:szenarien}

Um in dieser Arbeit etwas konkreter zu werden, wird ein Radarsystem ausgewählt und genauer betrachtet.
Da sonst alles nur sehr vage Umschreibungen ohne konkretes Ziel wären. 
Dieses Kapitel definiert Szenarien und Anforderungen für ein ausgewähltes Radarsystem.

\section{Verschiedene Radarsysteme}

In Flugzeugen sind meist mehrere Radare für unterschiedliche Zwecke verbaut. 
Im Folgenden werden exemplarisch vier Systeme betrachtet. Gezeigt werden jeweils Radartyp, 
typische Einsatzprofile und eine kurze Spezifikationsliste.

Hinweis zu Quellenangaben: 
\textbf{(N)} = Daten aus nicht belastbaren Quellen (z.\,B.\ Wikipedia oder Foren), 
\textbf{(T)} = typische Werte bekannt aus Vorlesungen wie Einführung autonomes Fahren.

\subsection*{MiG\textendash25 (Smerch--A, Kampfjet)\footnote{\citetitle{MiGAlley_Smerch} \parencite{MiGAlley_Smerch}}}
\begin{itemize}
  \item \textbf{Rolle/Typ:} Bordradar zur Luftzielsuche und \-verfolgung
  \item \textbf{Bezeichnung:} Smerch\textendash A
  \item \textbf{Frequenz/Band/Wellenlänge:} 10\,GHz (3\,cm); Variante A1: 15\,GHz (2\,cm)
  \item \textbf{Sendeleistung (Spitze):} \textbf{(N)} 600\,kW
  \item \textbf{Sendeleistung (Durchschnitt):} \textbf{(N)} 1–6\,kW
  \item \textbf{Wellenform/Puls:} \textbf{(N)} Pulsradar, Pulsbreite 0{,}025\,\textmu s
  \item \textbf{Reichweite/Detektion:} Bomber \(\approx\) 100\,km; Tracking \(\approx\) 50\,km (A1)
  \item \textbf{Scan/Abdeckung:} Fortlaufende Verbesserungen in Störfestigkeit und Tiefflug\-Clutter\-toleranz ab Serienstand
  \item \textbf{Antenne (Typ/Größe/Gain):} ---
  \item \textbf{Gewicht (System):} ca. 500\,kg
  \item \textbf{Abmessungen (H\(\times\)B\(\times\)T)/Volumen:} ---
  \item \textbf{Hinweise:} \textbf{(N)} keine Look\-down/Shoot\-down\-Fähigkeit
\end{itemize}

\subsection*{F\textendash14 (AWG\textendash9, Kampfjet)\footnote{\citetitle{Forecast1999_AWG9_APG71} \parencite{Forecast1999_AWG9_APG71}}}
\begin{itemize}
  \item \textbf{Rolle/Typ:} Bordradar
  \item \textbf{Bezeichnung:} AWG\-9
  \item \textbf{Frequenz/Band/Wellenlänge:} 8–12\,GHz (X\-Band)
  \item \textbf{Sendeleistung (Spitze):} 10\,kW
  \item \textbf{Sendeleistung (Durchschnitt):} 7\,kW (Pulse\-Doppler), 0{,}5\,kW (Pulsmodus)
  \item \textbf{Wellenform/Puls:} Pulsbreiten: Puls 0{,}4/50\,\textmu s; Pulse\-Doppler 0{,}4/1{,}3/2{,}0/2{,}7\,\textmu s
  \item \textbf{Reichweite/Detektion:} maximale Suche 213\,km
  \item \textbf{Scan/Abdeckung:} 80\(^\circ\)/s (horizontal); 2 Scans/s (vertikal)
  \item \textbf{Antenne (Typ/Größe/Gain):} \textbf{(N)} Slotted planar array; 14{,}2\,cm
  \item \textbf{Gewicht (System):} 590\,kg
  \item \textbf{Abmessungen (H\(\times\)B\(\times\)T)/Volumen:} 0{,}79\,m\(^3\)
  \item \textbf{Hinweise:} \textbf{(N)} Zielhöhenbereich ca. 50\,ft bis 80{,}000\,ft; Zielgeschwindigkeiten von niedriger Unterschall bis \(>\) Mach\,3
\end{itemize}

\subsection*{A\textendash320 (Honeywell RDR\textendash4000, Passagierflugzeug)\footnote{\citetitle{Honeywell_RDR4000_Specs} \parencite{Honeywell_RDR4000_Specs}}}
\begin{itemize}
  \item \textbf{Rolle/Typ:} Wetterradar
  \item \textbf{Bezeichnung:} Honeywell RDR\-4000
  \item \textbf{Frequenz/Band/Wellenlänge:} Senderfrequenz 9{,}375\,GHz
  \item \textbf{Sendeleistung (Spitze):} \textbf{(N)} 900\,W, Pulsbreite 12\,\textmu s
  \item \textbf{Sendeleistung (Durchschnitt):} \textbf{(N)} 40\,W, Pulszüge 275\,\textmu s
  \item \textbf{Wellenform/Puls:} Pulse Compression
  \item \textbf{Reichweite/Detektion:} 593\,km (Weather \& Ground Map); 111\,km (Turbulence); 9\,km (Windshear)
  \item \textbf{Scan/Abdeckung:} Azimut \(\pm 80^\circ\) (Weather/Ground Map), \(\pm 40^\circ\) (Windshear); Scanrate bis 90\(^\circ\)/s
  \item \textbf{Antenne (Typ/Größe/Gain):} Flat\-Plate\-Optionen 30"/24"/18"; Gain 34{,}8/33{,}0/31{,}0\,dBi; 3\-dB\-Breite 3{,}0/4{,}2/5{,}6\(^\circ\)
  \item \textbf{Gewicht (System):} Processor 4{,}76\,kg; Transmitter/Receiver 2{,}31\,kg; Antenne 7{,}3\,kg (Single), 13{,}4\,kg (Dual)
  \item \textbf{Abmessungen (H\(\times\)B\(\times\)T)/Volumen:} ---
  \item \textbf{Hinweise:} Umgebung \(-55^\circ\)C bis \(+70^\circ\)C; Software nach RTCA DO\-178B Level C; Minimum Discernible Signal \(-124\)\,dBm
\end{itemize}

\subsection*{Referenz (Automotive, Frontradar)\footnote{\citetitle{Bosch_FrontRadar_2025} \parencite{Bosch_FrontRadar_2025}}}
\begin{itemize}
  \item \textbf{Rolle/Typ:} Abstandsmessung
  \item \textbf{Bezeichnung:} Frontradar
  \item \textbf{Frequenz/Band/Wellenlänge:} 76–81\,GHz
  \item \textbf{Sendeleistung (Spitze):} ---
  \item \textbf{Sendeleistung (Durchschnitt):} \textbf{(T)} 0{,}1–1\,W
  \item \textbf{Wellenform/Puls:} \textbf{(T)} FMCW, Chirp\-Sequenz ca. 10\,\textmu s pro Rampe
  \item \textbf{Reichweite/Detektion:} bis zu 530\,m
  \item \textbf{Scan/Abdeckung:} \textbf{(N)} horizontal \(\pm 38^\circ\)
  \item \textbf{Antenne (Typ/Größe/Gain):} \textbf{(N)} 10 Elemente, jeweils \(d=3{,}2\) mm
  \item \textbf{Gewicht (System):} ---
  \item \textbf{Abmessungen (H\(\times\)B\(\times\)T)/Volumen:} 56\,mm \(\times\) 76\,mm \(\times\) 20\,mm
  \item \textbf{Hinweise:} Gesamtverbrauch < 3{,}5\,W; optional 5 Sensoren im Fahrzeug
\end{itemize}



\section{Auswahl des Radarsystems}
Da das System beispielhaft dienen soll, sind Einsatzzweck oder Art des Systems zweitrangig. 
Die Objektklassifikation wird für diese Arbeit aufgrund dessen gewählt, dass sie verhältnismäßig verständlich 
und nachvollziehbar ist. In dieser Arbeit wird beispielhaft eine Onboard-Wetterphänomensklassifikation als 
Funktion ausgewählt. Wetterphänomene zu identifizieren bietet hohes Potential, da aufgrund der Komplexität 
klassische Methoden schnell an ihre Grenzen geraten. Und zusätzliche Informationen über das Wetter 
die Crew bei Entscheidungen wie der Routenplanung unterstützen kann.


\section{Wetterklassifikation}

Ein Bordwetterradar dient der frühzeitigen Erkennung konvektiver Wetterlagen, 
um die Flugbahn vorausschauend zu planen und gefährliche Bereiche zu meiden. 
Die folgenden Grundlagen stützen sich überwiegend auf das Airbus-Magazin 
\enquote{Safety First}\footnote{\citetitle{marconnet2016optimumWeatherRadar} \parencite{marconnet2016optimumWeatherRadar}}.

\subsection*{Grundlagen und Physik}
Die drei Hauptgefahren sind Turbulenzen, Hagel und Windscherungen. Die Anzeige 
beruht auf der Reflexion flüssigen Wassers. Eis, trockener Hagel oder Schnee 
reflektieren schwächer und können trotz realer Gefahr unauffällig erscheinen. 
Die Stärke eines Echos entspricht daher nicht in jedem Fall der tatsächlichen Gefährdung. 
Das Echo hängt von Tropfengröße, Zusammensetzung und Menge ab. Ein Gewitter zeigt je 
nach Höhe unterschiedliche Reflektivität, da die Konzentration flüssigen Wassers mit der Höhe abnimmt.

\subsection*{Grenzen und typische Fehlinterpretationen}
Starke Niederschläge können Abschattungen verursachen, sodass scheinbar leere 
Bereiche hinter intensiven Echos entstehen. Turbulenz ohne Niederschlag wird nicht erfasst. 
Die Turbulenzanzeige bezieht sich auf niederschlagsgebundene Bereiche (Dopplereffekt). 
Form- und Strukturhinweise (z.\,B. ausgeprägte Kerne mit schwächeren Bereichen darüber) 
deuten auf starke Auf- und Abwinde hin. Schwach reflektierende Wolken über Land können auf starke Turbulenzen 
hinweisen, während nasse Wolken über dem Meer typischerweise weniger riskant sind.

\subsection*{Bedienhinweise und Scan-Empfehlungen}
Empfohlen wird, unterhalb der Gefriergrenze zu scannen und anschließend durch vertikale 
Scans sowie die Reflektivität aus höheren Regionen weitere Rückschlüsse zu ziehen. 
Turbulenzen sind nicht auf das Innere von Wolken begrenzt. Die Anwesenheit von Hagel 
variiert mit der Höhenlage. Hagel außerhalb von Wolken deutet auf starke Abwinde hin, 
feuchte Luft dagegen auf eine eher geringe Hagelwahrscheinlichkeit.

\subsection*{Erkennungsmerkmale}
Turbulenzen können anhand von Wolkenformen und häufigen Blitzen erkannt werden. Bestimmte 
Wolkenformen lassen auf schweren Hagel und starke Vertikalwinde schließen. Schnelle 
Formänderungen weisen auf hohe Wetteraktivität hin.

\subsection*{Taktische Ableitung und Abstände}
Grundsatz ist das weiträumige Meiden konvektiver Zellen, bevorzugt seitlich statt vertikal. 
Gefahren konvektiven Wetters sollten innerhalb von mindestens 40\,NM erkannt werden; Anpassungen 
im Flug erfolgen innerhalb von 80\,NM, und im Reiseflug wird eine Wettervermeidung innerhalb 
von 160\,NM angestrebt. Niemals unterhalb einer konvektiven Wolke fliegen, da dort sehr starke 
Turbulenzen auftreten können. Ein Überflug kommt nur in Betracht, wenn ein Ausweichen nicht möglich ist.










% -------------------------
% NEW CHAPTER: Ausgewählte KI-Modelle
% -------------------------
\chapter{Ausgewählte KI-Modelle}\label{chap:modelle}

Dieses Kapitel dient dazu einzelne relevante KI Modelle besser zu verstehen und eines davon 
auszuwählen, um dieses im Kontext der gewählten Wetterklassifikation weiter zu untersuchen.


\section{Modell-Steckbriefe}

\subsection*{CNN (Convolutional Neural Networks)} 
\begin{itemize}
  \item \textbf{Zweck:} Pixel/Region-Segmentierung \& Detektion.
  \item \textbf{Typische Variante:} YOLO
  \item \textbf{Input/Features:} 2D Radardaten (z.B. Winkel, Distanz), Radarrohdaten möglich.
  \item \textbf{Stärken:} Hohe Genauigkeit; lernt komplexe räumliche Muster; edge-tauglich mit  Quantisierung.
  \item \textbf{Schwächen:} Daten-/Label-intensiv; geringere Erklärbarkeit.
  \item \textbf{Latenz/Compute:} 10--100\,ms
  \item \textbf{Modellgröße:} ca. 2.000.000 Parameter.
\end{itemize}

\subsection*{Baumbasierte Ensembles (GBDT)}
\begin{itemize}
  \item \textbf{Zweck:} Tabellarische Klassifikation auf Feature-Vektoren.
  \item \textbf{Typische Variante:} Random Forest
  \item \textbf{Input/Features:} Texturen, SNR/Geometrie, Kontext (Temperatur/Höhe).
  \item \textbf{Stärken:} Schnell, robust bei kleinen/mittelgroßen Datensätzen; gut erklärbar; geringe Latenz.
  \item \textbf{Schwächen:} Räumliche Nachbarschaft nur indirekt; ungeeignet für rohe Bilder/Sequenzen.
  \item \textbf{Latenz/Compute:} $< 5$\,ms/Objekt; Echtzeitfähig
  \item \textbf{Modellgröße:} ca. 10.000 Knoten/Blätter.
\end{itemize}

\subsection*{Vision Transformer}

\begin{itemize}
  \item \textbf{Zweck:} Patch-/Bildbasierte Klassifikation, Segmentierung und Detektion.
  \item \textbf{Typische Variante:} ViT-B.
  \item \textbf{Input/Features:} Mehrkanal-Radar-Slices (z.\,B.\ Reflektivität, Spektralbreite, \text{SNR}).
  \item \textbf{Stärken:} Globaler Kontext durch Self-Attention; durch Tokenisierung auflösungsrobust.
  \item \textbf{Schwächen:} Hoher Daten- und Compute-Bedarf; Sensitivität gegenüber Patch-/Token-Design.
  \item \textbf{Latenz/Compute:}  $\sim$10--40\,ms pro 512$\times$512 (auf einer GPU).
  \item \textbf{Modellgröße:} ca. 20.000.000 Parameter.
\end{itemize}


\subsection*{Clustering}
\begin{itemize}
  \item \textbf{Zweck:} Objektbildung, Vor-/Nachverarbeitung, Tracking.
  \item \textbf{Typische Variante:} DBSCAN
  \item \textbf{Input/Features:} (Reflektivität, Textur) und Koordinaten; für Track-Kontext.
  \item \textbf{Stärken:} Formflexibel, label-arm; stabilisiert Objekte/Tracks und verbessert nachgelagerte Klassifikationen.
  \item \textbf{Schwächen:} Parameter-sensitiv; keine semantischen Labels.
  \item \textbf{Latenz:} Optional echtzeitfähig
  \item \textbf{Modellgröße:} Feature anhängig.
\end{itemize}

\section{Modellvergleich und Auswahl}

\small
\begin{longtable}{p{0.16\linewidth}p{0.17\linewidth}p{0.16\linewidth}p{0.18\linewidth}p{0.2\linewidth}}
\caption{Vergleich der Modelle mit Anforderungen}\\
\toprule
\textbf{Kriterium} &
\textbf{CNN} &
\textbf{Baumbasiert} &
\textbf{Transformer} &
\textbf{Clustering} \\
\midrule
\endfirsthead

\toprule
\textbf{Kriterium} &
\textbf{CNN} &
\textbf{Baumbasiert} &
\textbf{Transformer} &
\textbf{Clustering} \\
\midrule
\endhead

\midrule
\multicolumn{5}{r}{Fortsetzung auf der nächsten Seite}\\
\midrule
\endfoot

\bottomrule
\endlastfoot

\textbf{Passung zu Szenarien} &
\textbf{Hoch} &
\textbf{Hoch} mit guten Features &
\textbf{Hoch} auch großräumige Muster &
\textbf{Niedrig als Primär} formt Objekte, aber keine Klassen \\

\textbf{Leistung (F1, $P_d$/FAR)} &
\textbf{Hoch} bei genügend Daten &
\textbf{Hoch} mit guten Features &
\textbf{Hoch} bei genügend Daten &
\textbf{Parameterabhängig} \\

\textbf{Latenz} &
\textbf{Erfüllbar}&
\textbf{Sehr gut} &
\textbf{mittel} &
\textbf{Gut} \\

\textbf{Datenbedarf/ Labeling} &
\textbf{Hoch}  &
\textbf{Mittel} &
\textbf{Hoch} &
\textbf{Gering} \\

\textbf{Erklärbarkeit} &
\textbf{Mittel}  &
\textbf{Gut}  &
\textbf{Schlecht}  &
\textbf{Schlecht} keine Labels \\

\textbf{Harware- anforderung} &
\textbf{Mittel}  &
\textbf{Niedrig} &
\textbf{Hoch}  &
\textbf{Niedrig} \\

\textbf{Alleinstellungs-Tauglichkeit} &
\textbf{Ja}  &
\textbf{Ja} bei guten Features &
\textbf{Ja}  &
\textbf{Nein} Vor-/ Nachverarbeitung \\
\end{longtable}

\medskip
Als alleinstehendes Modell empfehlen sich Baummodelle (sofern man gute Features hat), als typisches Modell wird Random Forest verwendet. 
Diese sind alleinstellungsfähig, was für diese Arbeit wichtig ist, um zusätzliche Komplexität zu vermeiden. 
Zudem sind sie echtzeitfähig, gut erklärbar und haben eine relativ geringe Anforderung an die Leistung der Hardware. 
Das hat den Vorteil, dass es nicht bereits an der Hardwarezulassung scheitert.

\section{Erklärung eines Random Forest}
Ein Random Forest ist ein Ensemble vieler Entscheidungsbäume. Jeder Baum wird aus einer zufälligen Stichprobe 
der Trainingsdaten aufgebaut und betrachtet an jedem Knoten nur eine zufällige Auswahl der Merkmale. Der Baum 
stellt einfache Schwellenfragen, teilt die Daten Schritt für Schritt und gibt am Blatt eine Klassenentscheidung 
(oder einen Zahlenwert) aus. Für neue Daten werden die Ergebnisse aller Bäume zusammengeführt: Bei Klassifikation 
zählt die Mehrheit, bei Regression der Mittelwert. Durch die zufällige Daten- und Merkmalswahl werden die Bäume 
voneinander unabhängiger, die Streuung sinkt, und die Gesamtvorhersage wird robuster.

\begin{figure}[h!]
\centering
\begin{tikzpicture}[
  level 1/.style={sibling distance=46mm, level distance=14mm},
  level 2/.style={sibling distance=28mm, level distance=14mm},
  edge from parent/.style={draw,-Stealth},
  every node/.style={font=\small},
  decision/.style={rectangle,draw,rounded corners,align=center,inner sep=2pt,minimum width=22mm},
  leaf/.style={ellipse,draw,fill=gray!10,align=center,inner sep=2pt,minimum width=16mm}
]
\node[decision]{\textbf{Frage 1}\\$f_1 \le \theta_1$?}
  child { node[decision]{\textbf{Frage 2}\\$f_2 \le \theta_2$?}
    child { node[leaf]{\textbf{Blatt}\\Klasse A} edge from parent node[left]{Ja} }
    child { node[leaf]{\textbf{Blatt}\\Klasse B} edge from parent node[right]{Nein} }
    edge from parent node[left]{Ja}
  }
  child { node[decision]{\textbf{Frage 3}\\$f_3 \le \theta_3$?}
    child { node[leaf]{\textbf{Blatt}\\Klasse B} edge from parent node[left]{Ja} }
    child { node[leaf]{\textbf{Blatt}\\Klasse C} edge from parent node[right]{Nein} }
    edge from parent node[right]{Nein}
  };
\end{tikzpicture}
\caption{Einfache Skizze eines Entscheidungsbaums: An jedem Knoten wird eine Schwellenfrage gestellt; am Blatt steht die Klassenentscheidung.}
\end{figure}

\subsection*{Random Forest im Sicherheitskontext}

Random Forest eignet sich für sicherheitsrelevante Anwendungen, da das Verfahren von Natur aus 
robust gegen Overfitting ist. Die baumstrukturierten Regeln fördern die Erklärbarkeit: 
Entscheidungen lassen sich auf konkrete Feature-Schwellen und Pfade zurückführen. 
Gleichzeitig ist die Nachverfolgbarkeit sicherheitskritischer Unterscheidungen gegeben, 
weil Randfälle und knappe Entscheidungsmargen über Blattverteilungen, Knoten beziehungsweise 
Schwellen sowie fallbezogene Pfade identifiziert und auditiert werden können; kritische Fälle sind dadurch transparent analysierbar.

Zu weiteren Vorteilen zählt eine solide Basisleistung mit wenig Tuning, 
zudem ist das Modell über Feature-Wichtigkeiten gut erklärbar. Auf der anderen Seite ist Random Forest 
für eine Pixel-Segmentierung wenig geeignet. Bei vielen irrelevanten Features können tiefe Bäume Rauschen 
aufsammeln. Ohne geeignete Klassengewichte oder Sampling-Strategien können zudem Minderheitsklassen leicht benachteiligt werden.


\section{Einbindungsmöglichkeiten}
Ein Random Forest lässt sich entlang des gesamten Lebenszyklus eines Radarsystems nutzen. 
Er kann von der datengetriebenen Analyse in der Entwicklung bis zur begleitenden Entscheidungsunterstützung an Bord eingesetzt werden. 
Die folgenden Beispiele zeigen mögliche Integrationspunkte.

\medskip
\begin{itemize}
  \item \textbf{Musteridentifikation \& Datendarstellung in der Entwicklung}\\
  Wiederkehrende Strukturen in Radardaten können systematisch erkannt und als
  verständliche Merkmalsbilder aufbereitet werden. 
  Daten können visualisiert und ein konsistenten Musterkatalog 
  der später als Grundlage für die Klassifikation überführt werden.

  \item \textbf{Assistierendes System neben dem klassischen Radar}\\
  Das Modell läuft parallel und liefert zusätzliche Informationen zur Wetterlage
  z.\,B.\ Klassenhinweise oder Trends. Die Ausgabe ergänzt die bestehende Anzeige,
  ohne deren Logik zu ersetzen und unterstützt die Crew mit klaren, kompakten Hinweisen für die
  taktische Entscheidungsfindung.

  \item \textbf{Ersatz/Ergänzung klassischer Schwellen-Entscheidungsbäume}\\
  Anstelle fest verdrahteter Schwellwerte werden datengetriebene Entscheidungsgrenzen verwendet,
  die sich aus gelabelten Beispielen ableiten lassen. Das reduziert manuelles Tuning und 
  ermöglicht feinere Trennungen zwischen ähnlichen Phänomenen.

  \item \textbf{Kombination mit Kalman-Filter}\\
  Das Modell ersetzt die Zustandsschätzung klassischer Filter nicht, sondern füttert sie:
  Es liefert stabilisierte Beobachtungen und Klassifikationsergebnisse,
  die anschließend von einem Schätzer weiterverarbeitet werden. So verbleibt die Bahnverfolgung im
  Filter, während die inhaltliche Deutung der Echos und ihre zeitliche Glättung modellbasiert verbessert werden.
  Der Kalman-Filter wurde gewählt, da dessen Output
  probabilistisch ist und durch die Ergänzung eines ebenfalls probabilistischen Algorithmus nicht
  zwangsläufig in seiner Stabilität beeinträchtigt wird.
\end{itemize}

\medskip
Für diese Arbeit steht der Einsatz an Bord als eigenständige Assistenz im Fokus. 
Eine enge Kopplung oder gar ein Ersatz etablierter Software würde zusätzliche Komplexität 
und Zertifizierungsaufwand erzeugen, ohne den Mehrwert in diesem Rahmen zu klären. 




\section{Mögliche Szenarien beim Einsatz}
Beim parallelen Betrieb des klassischen Verfahrens und des Random-Forest ergeben sich wiederkehrende Entscheidungssituationen.
Die folgenden Punkte skizzieren naheliegende Optionen.


\begin{table}[h]
\centering
\renewcommand{\arraystretch}{1.12}
\begin{tabular}{p{0.36\linewidth} p{0.58\linewidth}}
\toprule
\textbf{Situation} & \textbf{Empfehlung} \\
\midrule
ML-Pfad meldet ein Ziel, der klassische Pfad nicht &
Möglich sind Frühtreffer oder ein Fehlalarm. Eine vorsichtige Meldung und Bestätigung über mehrere Messungen ist sinnvoll. \\
\addlinespace[4pt]
Klassischer Pfad meldet ein Ziel, der ML-Pfad nicht &
Konservativ entscheiden und den Befund des klassischen Pfads bevorzugen. \\
\addlinespace[4pt]
Beide melden kein Ziel &
Es kann davon ausgegangen werden, dass kein Ziel erkennbar ist und der Normalbetrieb kann fortgeführt werden. \\
\addlinespace[4pt]
Beide melden ein Ziel und stimmen überein &
Der Befund kann übernommenwerden und Zusatzinformationen aus dem ML-Pfad können unterstützen. \\
\addlinespace[4pt]
Beide melden ein Ziel, aber mit abweichender Einstufung &
Je nach Gefahrenlage konservativ entscheiden, im Zweifel die höhere Einstufung wählen oder den klassischen Pfad bevorzugen. \\
\bottomrule
\end{tabular}
\caption{Paralleler Betrieb von klassischem Verfahren und ML-Pfad: Entscheidungssituationen und Handlungsempfehlungen}
\end{table}






% -------------------------
% Anforderungsfelder & Referenzrahmen
% -------------------------
\chapter{Anforderungen \& Zuordnung}\label{chap:anforderungen}
In diesem Kapitel werden Anforderungen an das Beispiel des Randomforest zur Onbord-Wetterklassifikation gestellt 
und den Anforderungen aus der Luftfahrt zugeordnet.


\section{Anforderungen Wetterklassifikation und Random Forest}
\label{sec:req-radar-ki}

Dieser Abschnitt beschreibt Anforderungen an ein Klassifikationssystem auf Basis eines Bordwetterradars in Kombination 
mit einem Random Forest. Die Anforderungen sind einsatzzweckbezogen, qualitativ formuliert 
und mit IDs versehen, um sie später mit der Abdeckung der ISO 8800 abzugleichen.

\subsection*{ODD — Geltungsbereich}
\textbf{R-KI-01 ODD definieren und prüfen:} \\
Betriebsmodi des Radars, relevante Höhen- und Reichweitenbänder sowie Scanverfahren müssen festgelegt und laufend geprüft werden; außerhalb dieses Rahmens ist die Ausgabe als unsicher zu kennzeichnen.

\subsection*{Robustheit}
\textbf{R-KI-02 Konservative Behandlung schwacher oder fehlender Echos:}\\
Bereiche mit geringer Reflektivität oder möglicher Abschattung müssen als potenziell risikobehaftet erkannt und konservativ bewertet werden.

\textbf{R-KI-03 Vertikalstruktur als Gefahrenhinweis:} \\
Aus Tilt- und Vertikalscans abgeleitete Muster wie starke Kerne unten und schwächere Echos darüber müssen als Hinweis auf starke Auf- und Abwinde erkannt werden.

\textbf{R-KI-04 Form- und Dynamikmuster:} \\
Charakteristische Zellformen sowie schnelle Formänderungen müssen als erhöhte Wetteraktivität erkannt und mit erhöhter Priorität ausgegeben werden.

\textbf{R-KI-05 Turbulenz richtig einordnen:} \\
Niederschlagsgebundene Turbulenz muss erkannt werden; Turbulenz ohne Niederschlag ist als nicht detektierbar zu kennzeichnen und darf nicht als Entwarnung ausgelegt werden.

\subsection*{OOD — Randfälle und Abweichungen}
\textbf{R-KI-06 Randfälle explizit behandeln:} \\
Konstellationen wie Hagel außerhalb von Wolken, schwach reflektierende Wolken über Land oder starke Abschattungen müssen als Randfälle markiert und konservativ bewertet werden.

\textbf{R-KI-07 Drift überwachen:} \\
Änderungen der Datenlage durch Jahreszeit, Region oder Flughöhe müssen als Verteilungsdrift erkennbar sein; bis zur Klärung ist die Ausgabe konservativ zu halten und der Geltungsbereich bei Bedarf einzuschränken.

\subsection*{Betrieb und HMI}
\textbf{R-KI-08 Ausgabe auf Meiden ausrichten:} \\
Die Anzeige muss klare, kompakte Hinweise für seitliches Ausweichen liefern und keine Durchflugentscheidung ableiten.

\textbf{R-KI-09 Erklärbarkeit sicherstellen:} \\
Entscheidungen müssen auf Feature-Schwellen und Entscheidungswege zurückführbar sein; beitragende Merkmale sind zur Auditierbarkeit zu protokollieren.

\textbf{R-KI-10 Reproduzierbarkeit gewährleisten:} \\
Gleiche Eingaben müssen bei gleicher Modell- und Datenversion zu gleichen Ausgaben führen; Modell-, Daten- und Parameterstände sind mitzuloggen.



% -------------------------
% -------------------------

\section{Zuordnung Random-Forest-Anforderungen zu Luftfahrt-Forderungen}

\newcolumntype{L}[1]{>{\raggedright\arraybackslash}p{#1}}

\begin{longtable}{L{0.13\linewidth}L{0.42\linewidth}L{0.22\linewidth}}
\caption{Mapping der RF-spezifischen Anforderungen (R-KI-01…10) auf die Luftfahrt-Forderungen (F-Avi-01…08).}\\
\toprule
\textbf{RF-Req ID} & \textbf{RF-Anforderung (Kurztext)} & \textbf{Zuordnung F-Avi (IDs)} \\
\midrule
\endfirsthead
\toprule
\textbf{RF-Req ID} & \textbf{RF-Anforderung (Kurztext)} & \textbf{Zuordnung F-Avi (IDs)} \\
\midrule
\endhead
\midrule
\multicolumn{3}{r}{Fortsetzung auf nächster Seite}\\
\midrule
\endfoot
\bottomrule
\endlastfoot

R-KI-01 & ODD definieren und prüfen (außerhalb ODD als unsicher kennzeichnen) & F-Avi-02 \\[2pt]

R-KI-02 & Konservative Behandlung schwacher/fehlender Echos (Abschattung berücksichtigen) & F-Avi-05, F-Avi-02 \\[2pt]

R-KI-03 & Vertikalstruktur als Gefahrenhinweis (Tilt-/Vertikalscans) & F-Avi-05 \\[2pt]

R-KI-04 & Form- und Dynamikmuster als erhöhte Aktivität priorisieren & F-Avi-05, F-Avi-08 \\[2pt]

R-KI-05 & Turbulenz richtig einordnen (nur niederschlagsgebunden detektierbar) & F-Avi-05, F-Avi-08 \\[2pt]

R-KI-06 & Randfälle explizit behandeln (Hagel außerhalb Wolken, schwache Landwolken, Abschattung) & F-Avi-03, F-Avi-02 \\[2pt]

R-KI-07 & Drift überwachen und konservativ reagieren & F-Avi-06, F-Avi-04 \\[2pt]

R-KI-08 & Ausgabe auf Meiden ausrichten (klare, kompakte Hinweise; keine Durchflugentscheidung) & F-Avi-01, F-Avi-08 \\[2pt]

R-KI-09 & Erklärbarkeit sicherstellen (Feature-Schwellen/Pfade protokollieren) & F-Avi-01 \\[2pt]

R-KI-10 & Reproduzierbarkeit gewährleisten (gleiche Eingabe → gleiche Ausgabe; Versionen loggen) & F-Avi-06 \\

\end{longtable}

\medskip
In Summe sind alle Random-Forest-Anforderungen den Luftfahrt-Forderungen zuordenbar, eine echte
„Nicht-Zuordnung“ ergibt sich nicht. Projektspezifische Präzisierungen sind zu berücksichtigen, in den allgemeinen 
Forderungen sind sie lediglich konzeptionell gehalten.







% -------------------------
% CHAPTER 7: Bewertung der Übertragbarkeit
% -------------------------
\chapter{Bewertung der Übertragbarkeit}\label{chap:uebertragbarkeit}
In diesem Kapitel geht es darum die Übertragbarkeit des gewählten Beispiels auf existierende Normen zu 
übertragen. Das in dieser Arbeit gewählte Beispiel ist ein Random Forest KI Modell, um 
Muster und Datensätze in der Entwicklung zu erkennen und darzustellen. Für die Klassifikation von  
Wetterphänomenen soll es unter den Vorraussetzungen
eines Onboard Radarsystems in einem Flugzeug betrachtet werden.



\section{Abgeleitete ISO~8800-Abdeckung (aus R-KI \texorpdfstring{$\rightarrow$}{→} F-Avi)}

\subsection*{Methode}
Die Abdeckung pro RF-Anforderung wird aus der Zuordnung R-KI \(\rightarrow\) F-Avi und der dort ausgewiesenen ISO-Abdeckung hergeleitet:
Sind alle zugeordneten F-Avi-Punkte „voll“, dann \covV; andernfalls \covT\ (konservativ).

\newcolumntype{L}[1]{>{\raggedright\arraybackslash}p{#1}}

\begin{longtable}{L{0.12\linewidth}L{0.36\linewidth}L{0.24\linewidth}L{0.08\linewidth}L{0.18\linewidth}}
\caption{Abgeleitete Abdeckung der ISO/PAS~8800 nach RF-Anforderungen (R-KI-01…10).}\\
\toprule
\textbf{R-KI} & \textbf{Kurztext} & \textbf{ISO~8800 (Kapitel, abgeleitet)} & \textbf{Abd.} & \textbf{Hinweis} \\
\midrule
\endfirsthead
\toprule
\textbf{R-KI} & \textbf{Kurztext} & \textbf{ISO~8800 (Kapitel, abgeleitet)} & \textbf{Abd.} & \textbf{Hinweis} \\
\midrule
\endhead
\midrule
\multicolumn{5}{r}{Fortsetzung auf nächster Seite}\\
\midrule
\endfoot
\bottomrule
\endlastfoot

R-KI-01 & ODD definieren und prüfen (außerhalb ODD als unsicher kennzeichnen) & 9, 12, 14 & \covV & F-Avi-02 \covV \\[2pt]

R-KI-02 & Konservative Behandlung schwacher/fehlender Echos (Abschattung berücksichtigen) & 9, 12, 14 & \covV & F-Avi-05 \covV; F-Avi-02 \covV \\[2pt]

R-KI-03 & Vertikalstruktur als Gefahrenhinweis (Tilt-/Vertikalscans) & 9, 12 & \covV & F-Avi-05 \covV \\[2pt]

R-KI-04 & Form- und Dynamikmuster als erhöhte Aktivität priorisieren & 6, 8, 9, 12, 14 & \covT & F-Avi-05 \covV; F-Avi-08 \covT \\[2pt]

R-KI-05 & Turbulenz richtig einordnen (nur niederschlagsgebunden detektierbar) & 6, 8, 9, 12, 14 & \covT & F-Avi-05 \covV; F-Avi-08 \covT \\[2pt]

R-KI-06 & Randfälle explizit behandeln (Hagel außerhalb Wolken, schwache Landwolken, Abschattung) & 8, 9, 11, 12, 14 & \covV & F-Avi-03 \covV; F-Avi-02 \covV \\[2pt]

R-KI-07 & Drift überwachen und konservativ reagieren & 7, 8, 14 & \covT & F-Avi-06 \covT; F-Avi-04 \covV \\[2pt]

R-KI-08 & Ausgabe auf Meiden ausrichten (klare Hinweise, keine Durchflugentscheidung) & 6, 7, 8, 12, 14 & \covT & F-Avi-01 \covT; F-Avi-08 \covT \\[2pt]

R-KI-09 & Erklärbarkeit sicherstellen (Feature-Schwellen/Pfade protokollieren) & 6, 7 & \covT & F-Avi-01 \covT \\[2pt]

R-KI-10 & Reproduzierbarkeit gewährleisten (Versionen loggen) & 7, 8, 14 & \covT & F-Avi-06 \covT \\

\end{longtable}


\subsection*{Abdeckungs–Zwischenfazit}
Die RF-Anforderungen werden durch die ISO 8800 weitgehend adressiert, 
mit \textbf{voller Abdeckung} für den ODD-Rahmen samt konservativem Verhalten 
(Abstain/Degradation), für Datenaspekte und Randfälle (Repräsentativität, 
Abschattung/Unsicherheitszonen) sowie Vorgehensempfohlene Informationsgewinne 
(für strukturelle Wetterhinweise aus Vertikal-/Tilt-Scans). 

\medskip
\textbf{Teilweise abgedeckt} bleiben anwendungsnahe Punkte wie Form-/Dynamikmuster, 
die spezifische Einordnung der Turbulenzanzeige, Drift-Überwachung im Betrieb, HMI-gerechte 
Ausgabe (Meiden statt Durchflug), Erklärbarkeit (Feature-Schwellen/Pfade) und Reproduzierbarkeit 
über Versionen. Hier sind projektspezifische Präzisierungen nötig, um den prinzipienbasierten Rahmen 
der Norm in konkrete, radar- und einsatzbezogene Nachweise zu überführen.


\section{Untersuchung der teilweise abgedeckten Punkte}

Im Folgenden werden die in der Abdeckungstabelle als \covT gekennzeichneten RF-Anforderungen kurz diskutiert. 

\begin{itemize}
  \item \textbf{R-KI-04 Form- und Dynamikmuster priorisieren:}\\
  Gut erfüllbar. Wenn in den Merkmalen beispielsweise Formänderung (Reflektivität- und Ausdehnungs-Unterschiede über die Zeit) enthalten sind.

  \item \textbf{R-KI-05 Turbulenz richtig einordnen:}\\
  Mittelmäßig erfüllbar. Nicht erkennbare Turbulenzen können zumindest regelbasiert als nicht detektierbar gekennzeichnet werden.

  \item \textbf{R-KI-07 Drift überwachen und konservativ reagieren:}\\
  Mittelmäßig erfüllbar. Konsevative Schwellen können zwar gesetzt werden, aber Drift muss zusätzlich geprüft werden.

  \item \textbf{R-KI-08 Ausgabe auf Meiden ausrichten:}\\
  Gut erfüllbar. Random Forest Ausgaben können auf risikoorientierte Klassen übertragen werden.

  \item \textbf{R-KI-09 Erklärbarkeit sicherstellen:}\\
  Sehr gut erfüllbar. Entlang der Pfade beziehungsweise Knoten können Unterscheidungen nachverfolgt und auditiert werden.

  \item \textbf{R-KI-10 Reproduzierbarkeit gewährleisten:}\\
  Sehr gut erfüllbar. Ein Random Forest nimmt erfahrungsgemäß bei gleichen Umgebungen und Inputs die gleiche Struktur beziehungsweise Knoten und Unterscheidungen an.
\end{itemize}

\section{KO-Kriterien}
Die folgenden Punkte sind Mindestvoraussetzungen, damit die Funktion grundsätzlich für eine Zulassung
als Bordfunktion in Betracht kommt. Sie adressieren den Onboard-Einsatz, ein reiner Entwicklungsbetrieb
(ohne Bordfunktion) fällt nicht darunter.

\medskip
\begin{itemize}
  \item \textbf{AI-Level \& Autorität (EASA L1/L2):}
        Unterstützendes Modell für den Menschen, welches jederzeit übersteuerbar bleibt.\\
        Ist als Bordassistenzsystem entsprechend gewählt worden und somit umsetzbar.

  \item \textbf{DAL Level C/D:}
        Empfehlungsgemäß soll keine KI in hoch Risiko Bereichen verwendet werden (DAL A/B). \\
        Die Zielgrößen gilt es auszuarbeiten, sollte aber mit verhältnismäßigen Aufwand möglich sein. 
        Eine Wetterklassifikation entspricht typischerweise DAL C und liegt somit in einem realistischen Bereich.

  \item \textbf{Kein Online-Lernen im Betrieb:}
        Eingefrorenes Modell: Das Modell wird offline trainiert und zur Freigabe „gefroren“;
        Laufzeitanpassungen sind ausgeschlossen.\\ 
        Für einen Random Forest ist dies unproblematisch umsetzbar.

  \item \textbf{Kein Nachweis-Rabatt:}
        ML-Erkenntnisse ersetzen keine klassischen Nachweise: Jede abgeleitete Regel/Muster wird
        mit etablierten Verfahren verifiziert und in den Safety Case integriert (vollständige Nachweiskette). \\
        Dies gilt es auszuarbeiten, sollte aber mit verhältnismäßigen Aufwand möglich sein.

  \item \textbf{Zugelassene Hardware und Werkzeuge:}
        Ohne entsprechend zugelassene Hardware und Werkzeuge ist die Erbringung eines belastbaren 
        Sicherheitsnachweises erheblich erschwert.\\
        Ein Random Forest sollte mit derzeit zugelassener Hardware lauffähig sein. 
        Zugelassene Werkzeuge oder Nachweise für verwendete Werzeuge sollte in absehbarer Zeit beziehungsweise Aufwand möglich sein.
\end{itemize}


\section{Fazit der Übertragbarkeit}
Die Zuordnung zeigt, dass die für den betrachteten Anforderungen 
durch ISO 8800 mindestens teilweise und in der Praxis gut erfüllbar abgedeckt werden. 
Mit Blick auf eine Bordfunktion sind die KO-Kriterien realistisch erfüllbar. 
Ein Random Forest eignet sich dafür: rechnerisch leicht, deterministisch, erklärbar und mit etablierten Evidenzformen. 
Die noch teilabgedeckten Punkte sind projektspezifisch zu präzisieren, stellen jedoch keine prinzipiellen 
Hürden dar.

Den Honeywell-Dokumenten zum A320-Wetterradar zufolge fällt die Software nach DO-178B 
(Vorgänger von DO-178C) in DAL~C und kann nach EASA-Empfehlungen Anwendung finden. Insgesamt ist die Übertragbarkeit 
unter den definierten Einsatzgrenzen (L1, ODD, kein Online-Lernen) und mit den skizzierten Evidenzen ausreichend für 
eine zulassungsfähige Assistenzfunktion. Offene Punkte betreffen primär die projektspezifische Ausgestaltung, nicht 
die grundsätzliche Machbarkeit.
